% GENERATED FILE. Edit canonical lessons, then run npm run generate:books.
% canonical-source-hash: fnv1a64:20135390ed595042
% canonical-lessons: GE-C304-trauer, GE-C304-trauern, GE-C304-beileid, GE-C304-witwe, GE-C304-erben

\chapter{Die Trauer}
\label{ch:ge-die-trauer}
\hlchaptermodality{304}

\begin{chapteropening}
\textbf{By the end of this chapter:} \emph{I can talk about mourning: grief, to mourn, condolences, a widow and to inherit.}

Complete the last lesson of chapter 304: I can talk about mourning: grief, to mourn, condolences, a widow and to inherit.
\end{chapteropening}

\section[die Trauer]{die Trauer — the grief, mourning}

\label{lesson:GE-C304-trauer}



\begin{quote}
Before the new one: say the German for the coffin, then the German for the cemetery.
\end{quote}

\subsection*{You'll want to know: die Trauer}
\textbf{die Trauer} — \textquotedblleft{}the grief, mourning\textquotedblright{}.

\begin{quote}
Die Trauer in der Familie war groß.
\end{quote}

\begin{quote}
\emph{The grief in the family was great.}
\end{quote}

No plural. \emph{Traurig}, already known, is related.

\begin{cousinweb}
\emph{Trauer} comes from Middle High German \emph{trūre} \textquotedblleft{}mourning\textquotedblright{}, from Old High German \emph{trūrēn} \textquotedblleft{}to mourn\textquotedblright{} (EtymDB).
\end{cousinweb}

\begin{grammarlens}[title={What you've built}]
One more word for telling, explaining and describing at B1.
\end{grammarlens}

\subsection*{Your turn}
\begin{itemize}
\raggedright
  \item \emph{Say it:} \emph{die Trauer}
  \item \emph{Say it:} \emph{die Trauer}, once more
  \item \emph{Say it:} a sentence of your own with \emph{die Trauer}
  \item \emph{Recall:} say the German for the coffin, then the German for the cemetery, then say \emph{die Trauer} again
\end{itemize}

\begin{tcolorbox}[breakable,colback=teal!4,colframe=teal!35!black,title={Before you move on}]
What does die Trauer mean? (\textquotedblleft{}The grief, mourning\textquotedblright{}.) Say it once more.
\end{tcolorbox}
\section[trauern]{trauern — to mourn, to grieve}

\label{lesson:GE-C304-trauern}



\begin{quote}
Before the new one: say the German for the cemetery, then the German for the grief.
\end{quote}

\subsection*{You'll want to know: trauern}
\textbf{trauern} — \textquotedblleft{}to mourn, to grieve\textquotedblright{}.

\begin{quote}
Das ganze Dorf trauerte um den alten Arzt.
\end{quote}

\begin{quote}
\emph{The whole village mourned the old doctor.}
\end{quote}

\emph{Um jemanden trauern} — to mourn someone.

\begin{cousinweb}
\emph{trauern} is inherited from Middle High German \emph{trūren}, Old High German \emph{trūrēn} \textquotedblleft{}to mourn\textquotedblright{} (EtymDB).
\end{cousinweb}

\begin{grammarlens}[title={What you've built}]
One more word for telling, explaining and describing at B1.
\end{grammarlens}

\subsection*{Your turn}
\begin{itemize}
\raggedright
  \item \emph{Say it:} \emph{trauern}
  \item \emph{Say it:} \emph{trauern}, once more
  \item \emph{Say it:} a sentence of your own with \emph{trauern}
  \item \emph{Recall:} say the German for the cemetery, then the German for the grief, then say \emph{trauern} again
\end{itemize}

\begin{tcolorbox}[breakable,colback=teal!4,colframe=teal!35!black,title={Before you move on}]
What does trauern mean? (\textquotedblleft{}To mourn, to grieve\textquotedblright{}.) Say it once more.
\end{tcolorbox}
\section[das Beileid]{das Beileid — the condolences, sympathy}

\label{lesson:GE-C304-beileid}



\begin{quote}
Before the new one: say the German for the grief, then the German for to mourn.
\end{quote}

\subsection*{You'll want to know: das Beileid}
\textbf{das Beileid} — \textquotedblleft{}the condolences, sympathy\textquotedblright{}.

\begin{quote}
Wir sprachen der Familie unser Beileid aus.
\end{quote}

\begin{quote}
\emph{We offered our condolences to the family.}
\end{quote}

\emph{Herzliches Beileid} — my sincere condolences.

\begin{cousinweb}
\emph{Beileid} joins \emph{bei} \textquotedblleft{}with\textquotedblright{} and \emph{Leid} \textquotedblleft{}sorrow\textquotedblright{}, an inherited Germanic word whose English cousin is \emph{loath} (EtymDB): sorrow shared with someone.
\end{cousinweb}

\begin{grammarlens}[title={What you've built}]
One more word for telling, explaining and describing at B1.
\end{grammarlens}

\subsection*{Your turn}
\begin{itemize}
\raggedright
  \item \emph{Say it:} \emph{das Beileid}
  \item \emph{Say it:} \emph{das Beileid}, once more
  \item \emph{Say it:} a sentence of your own with \emph{das Beileid}
  \item \emph{Recall:} say the German for the grief, then the German for to mourn, then say \emph{das Beileid} again
\end{itemize}

\begin{tcolorbox}[breakable,colback=teal!4,colframe=teal!35!black,title={Before you move on}]
What does das Beileid mean? (\textquotedblleft{}The condolences, sympathy\textquotedblright{}.) Say it once more.
\end{tcolorbox}
\section[die Witwe]{die Witwe — the widow}

\label{lesson:GE-C304-witwe}



\begin{quote}
Before the new one: say the German for to mourn, then the German for the condolences.
\end{quote}

\subsection*{You'll want to know: die Witwe}
\textbf{die Witwe} — \textquotedblleft{}the widow\textquotedblright{}.

\begin{quote}
Die Witwe wohnte noch viele Jahre im Haus.
\end{quote}

\begin{quote}
\emph{The widow lived in the house for many more years.}
\end{quote}

Plural: \emph{die Witwen}. A man is \emph{der Witwer}.

\begin{cousinweb}
\emph{Witwe} is inherited from Old High German \emph{wituwa}, from Proto-Germanic \emph{widuwōn} \textquotedblleft{}widow\textquotedblright{} (EtymDB), the same source as English \emph{widow}.
\end{cousinweb}

\begin{grammarlens}[title={What you've built}]
One more word for telling, explaining and describing at B1.
\end{grammarlens}

\subsection*{Your turn}
\begin{itemize}
\raggedright
  \item \emph{Say it:} \emph{die Witwe}
  \item \emph{Say it:} \emph{die Witwe}, once more
  \item \emph{Say it:} a sentence of your own with \emph{die Witwe}
  \item \emph{Recall:} say the German for to mourn, then the German for the condolences, then say \emph{die Witwe} again
\end{itemize}

\begin{tcolorbox}[breakable,colback=teal!4,colframe=teal!35!black,title={Before you move on}]
What does die Witwe mean? (\textquotedblleft{}The widow\textquotedblright{}.) Say it once more.
\end{tcolorbox}
\section[erben]{erben — to inherit}

\label{lesson:GE-C304-erben}



\begin{quote}
Before the new one: say the German for the condolences, then the German for the widow.
\end{quote}

\subsection*{You'll want to know: erben}
\textbf{erben} — \textquotedblleft{}to inherit\textquotedblright{}.

\begin{quote}
Sie erbte das Haus ihrer Großmutter.
\end{quote}

\begin{quote}
\emph{She inherited her grandmother's house.}
\end{quote}

\emph{Das Erbe} — the inheritance.

\begin{cousinweb}
\emph{erben} is an inherited Germanic verb, which EtymDB traces to Proto-Germanic \emph{arbijaną} \textquotedblleft{}to inherit\textquotedblright{}.
\end{cousinweb}

\begin{grammarlens}[title={What you've built}]
One more word for telling, explaining and describing at B1.
\end{grammarlens}

\subsection*{Your turn}
\begin{itemize}
\raggedright
  \item \emph{Say it:} \emph{erben}
  \item \emph{Say it:} \emph{erben}, once more
  \item \emph{Say it:} a sentence of your own with \emph{erben}
  \item \emph{Recall:} say the German for the condolences, then the German for the widow, then say \emph{erben} again
\end{itemize}

\begin{tcolorbox}[breakable,colback=teal!4,colframe=teal!35!black,title={Before you move on}]
What does erben mean? (\textquotedblleft{}To inherit\textquotedblright{}.) Say it once more.
\end{tcolorbox}
